\chapter{The frustration-free area law by an exact sweep}
\label{ch:peps_area_law_frustration_free}

This chapter follows the reading note \cite{TNLean2026FrustrationFreeNote}, which
specializes the two-dimensional area law \cite[Theorem~1.1]{OpenAI2026AreaLaw}
to frustration-free Hamiltonians of finite range. Each local term annihilates
the ground vector, so a product of local filters suppresses the excited part of
any vector while a local perturbation of the ground vector spreads by an exactly
bounded distance. This replaces the time-filtered constraints and the
localization errors of the amplification stage.

\section{Contraction of one sweep}
\label{sec:ff_area_law_sweep}

\begin{lemma}[One sweep of projections]
    \label{thm:ff_area_law_sweep_contraction}
    \lean{ContinuousLinearMap.complementProd_apply_of_forall_apply_eq_zero, ContinuousLinearMap.inner_complementProd_of_forall_apply_eq_zero, ContinuousLinearMap.norm_sq_complementProd_le_of_gap, ContinuousLinearMap.norm_pow_complementProd_le_of_gap}
    \leanok
    Let $\Pi_1,\dots,\Pi_m$ be orthogonal projections on a Hilbert space, each
    failing to commute with at most $d$ of the others, and let $\Omega$ be a
    vector with $\Pi_i\Omega=0$ for all $i$. Suppose that $\gamma\ge0$ and
    $\langle v,\sum_i\Pi_iv\rangle\ge\gamma\|v\|^2$ for every $v\perp\Omega$.
    Put $D=(I-\Pi_m)\cdots(I-\Pi_1)$. Then $D\Omega=\Omega$,
    $\langle\Omega,Dv\rangle=\langle\Omega,v\rangle$ for every $v$, and for
    $v\perp\Omega$
    \[
        \bigl((d+1)^2+\gamma\bigr)\|Dv\|^2\le(d+1)^2\|v\|^2,\qquad
        \|D^tv\|\le\theta^t\|v\|,\qquad
        \theta=\bigl(1+\gamma/(d+1)^2\bigr)^{-1/2}.
    \]
    This is the contraction \cite[eq.~(sweep)]{TNLean2026FrustrationFreeNote}
    for filters that are projections; the note states it for positive
    contractions $k_i$ with the factors $(I-k_i)^{1/2}$. The restriction is
    recorded in \cite{gap:ffarealaw_sweep_range_projections}.
\end{lemma}

\begin{proof}\leanok
    The detectability lemma \cite{AnshuAradVidick2016} gives
    $\langle Dv,\sum_i\Pi_iDv\rangle\le d^2(\|v\|^2-\|Dv\|^2)$. Each factor
    $I-\Pi_i$ is self-adjoint and fixes $\Omega$, so $Dv\perp\Omega$ and the
    left side is at least $\gamma\|Dv\|^2$. Adding $(2d+1)\|Dv\|^2\le(2d+1)\|v\|^2$
    gives the first bound. Since $D$ preserves the orthogonal complement of
    $\Omega$, the second bound follows by induction on $t$.
\end{proof}

\begin{lemma}[Range projections of positive local terms]
    \label{thm:ff_area_law_range_projection}
    \lean{Matrix.PosSemidef.le_smul_supportProj, TNLean.PEPS.AreaLaw.FrustrationFree.supportProj_mem_supportedOperators}
    \leanok
    For a positive semidefinite matrix $A$, let $\Pi_A$ be the orthogonal
    projection onto the range of $A$, so that $\Pi_Av=0$ whenever $Av=0$. If
    $A\le J$, then $A\le J\Pi_A$. If $A$ acts on a set of sites $D$, so does
    $\Pi_A$.
\end{lemma}

\begin{proof}\leanok
    The projection $\Pi_A$ is $f(A)$ with $f(0)=0$ and $f(x)=1$ for $x\ne0$,
    and $x\le Jf(x)$ on the spectrum of $A$, which lies in $[0,J]$. An operator
    acts on $D$ exactly when it commutes with every operator acting on the
    complement of $D$, and functions of $A$ inherit this commutation.
\end{proof}
